% --- Start of corrected file for ANONYMOUS REVIEW ---
% !TEX program = pdflatex % Specify engine if needed by editor

% Use the IEEE Transactions on AI class (journal mode is default)
\documentclass[journal]{IEEEtran} 

% --- Essential Packages ---
\usepackage[utf8]{inputenc} % Input encoding
\usepackage[T1]{fontenc}    % Font encoding for better hyphenation
\usepackage{amsmath}
\usepackage{amssymb}        % Math symbols
\usepackage{amsfonts}       % Math fonts
\usepackage{bm}             % Bold math symbols (\bm)
\usepackage{graphicx}       % Include graphics
\usepackage{caption}        % Custom captions
\usepackage{microtype}      % Improved typography and justification
\usepackage{cite}           % Handles citations like [1]-[3]
\usepackage{soul}           % For \hl{} highlighting
\usepackage{siunitx}        % For \sisetup and S column type
\usepackage{multirow}       % Multi-row cells in tables
\usepackage{booktabs}       % Nicer tables
\usepackage{textcomp}       % Additional text symbols
% --- Optional/Situational Packages ---
\usepackage{tabularx}       
\usepackage[absolute,overlay]{textpos} 
\usepackage{algorithm}      
\usepackage{algpseudocode}  
\usepackage{enumitem}       
\usepackage[usenames,dvipsnames]{xcolor} 
\usepackage{hyperref}       
\usepackage{xcolor} % Required for defining colors
\usepackage{soul}   % Required for the \hl command

\sethlcolor{yellow}

% --- Document Settings ---
\captionsetup{font=footnotesize, labelfont=bf, skip=4pt} % IEEE style captions

% --- Title and Author (ANONYMIZED for double-blind review) ---
\title{Neural Contextual Anomaly Detection with Counterfactual Successor Memory in Space Systems}

% Placeholder for anonymous submission
\author{%
  \begin{tabular}{ccc}
    Andre Williams & Imadeldin Mahgoub & Hamid Akbarian\\
    \textit{Florida Atlantic University} & \textit{Florida Atlantic University} & \textit{Florida Atlantic University} \\
    \textit{Department of EECS} & \textit{Department of EECS} & \textit{Department of EECS} \\ 
    \textit{Boca Raton, US} & \textit{Boca Raton, US} & \textit{Boca Raton, US} \\
    \textit{andrewilliam2018@fau.edu} & \textit{mahgoubi@fau.edu} & \textit{hakbaria@fau.edu}
  \end{tabular}
}

% FOR DOUBLE-BLIND REVIEW: Running headers should be anonymized or removed.
\markboth{IEEE Transactions on Aerospace and Electronic Systems, SUBMISSION FOR REVIEW}%
{Submission}

% --- Document Start ---
\begin{document}

\maketitle

\begin{abstract}
Detecting sustained, contextual anomalies in spacecraft telemetry remains difficult because a long anomaly can contaminate the recent context used by contrastive, predictive, or reconstructive detectors. We introduce Neural Contextual Anomaly Detection with Counterfactual Successor Memory (NCAD-CS), a framework that asks a directly operational question: given the current context, what should the next telemetry segment have looked like under normal behavior? NCAD-CS stores normal context embeddings together with the normal successor segments that followed them during training. At inference time, nearest normal contexts retrieve plausible counterfactual successors, and the observed suspect segment is scored against this successor set. The resulting successor evidence is fused with local deviation evidence and a context-distance reliability signal, then converted into point-level events through a channel-adaptive threshold that combines the training successor distribution with an adaptive elbow-based score floor. On the public NASA SMAP/MSL telemetry benchmark, the accepted 15-epoch configuration evaluates 81 labeled channels with macro F1 of 0.3883, micro F1 of 0.4233, 34 channels above 0.5 F1, and 16 channels above 0.8 F1. The method is strong on several sustained and abrupt channels, including D-3, A-7, P-3, and T-5, but still fails on subtle or highly channel-specific regimes. Additional experiments with longer training, manifold-uncertainty scoring, and a multi-scale TCN encoder show that these variants do not yet improve the full benchmark as default replacements. The results support counterfactual successor retrieval as a promising direction while making clear that recall preservation and channel-adaptive event calibration remain central open problems.
\end{abstract}

\begin{IEEEkeywords}
Spacecraft Telemetry, Anomaly Detection, Contextual Anomaly, Sustained Anomaly, Temporal Convolutional Network, Counterfactual Successor Memory, Contrastive Learning, Adaptive Thresholding, F1 Score.
\end{IEEEkeywords}

\section{Introduction}
\label{sec:intro}

Monitoring spacecraft health through telemetry data is a cornerstone of modern space mission operations, essential for ensuring mission success, longevity, and safety \cite{meng2019spacecraft}. This data, a high-dimensional stream of measurements from thousands of onboard sensors, provides a continuous view into the operational state of complex subsystems. However, the sheer volume, velocity, and complexity of this data make manual analysis infeasible, necessitating the use of automated and intelligent anomaly detection systems. While conventional methods based on static, predefined operational limits have long been used, they are often insufficient for identifying the nuanced, context-dependent deviations that frequently precede significant system failures \cite{chandola2009anomaly}. The dynamic nature of space environments and the evolution of hardware behavior over a mission's lifetime demand more sophisticated approaches.

Recent advancements in machine learning, particularly in deep learning, have introduced powerful new paradigms for modeling complex temporal dependencies in time series data \cite{chalapathy2019deep}. Yet, several fundamental challenges persist in the space domain. These include the extreme rarity of anomalous events, which complicates model training; the high dimensionality of telemetry; and the critical need for interpretable and trustworthy results \cite{cuellar2024explainable}.

A particularly pernicious challenge is the detection of contextual and sustained anomalies. Unlike simple point anomalies (e.g., a sudden sensor spike), contextual anomalies are defined by their relationship to preceding data. A telemetry reading may be within its global operational limits, but anomalous given its immediate temporal context. Sustained anomalies are even more difficult, such as a gradual system drift or a shift into a new, off-nominal operational regime. Standard contrastive, reconstructive, and predictive models often fail in these scenarios due to \textit{context contamination}: as the anomaly persists, the model's immediate context window becomes filled with anomalous data, causing it to accept this new state as the "new normal" and suppress its anomaly score.

To address these specific challenges, this paper introduces Neural Contextual Anomaly Detection with Counterfactual Successor Memory (NCAD-CS). The central idea is to avoid asking only whether the current suspect window resembles its immediate context. Instead, NCAD-CS asks what successor segment should have followed the current context under normal behavior. A memory of normal context-successor pairs retrieves plausible counterfactual futures, and the observed suspect segment is evaluated against those retrieved normal futures. This reframes context reliability as a successor-consistency problem: even if the immediate context is ambiguous or partially contaminated, similar normal contexts can still provide a reference distribution for what should happen next.

This paper details the architecture and current evaluation status of NCAD-CS on the public NASA telemetry benchmark. Our experiments show that the framework can detect several sustained contextual anomalies well, particularly when the retrieved normal successors form a useful local reference set. They also reveal persistent weaknesses on subtle channels, flat or highly saturated regimes, and cases where high recall is obtained only by accepting many false positives. We therefore position NCAD-CS as a promising but still evolving framework built from robust multi-domain feature engineering, a hybrid-pooling Temporal Convolutional Network (TCN) encoder, counterfactual successor retrieval, local deviation evidence, and adaptive event calibration.

The remainder of this paper is organized as follows. Section~\ref{sec:related} reviews related work in spacecraft anomaly detection. Section~\ref{sec:dataset} describes the dataset used for evaluation. Section~\ref{sec:proposed_method} presents the NCAD-CS methodology. Section~\ref{sec:exp} reports the all-channel benchmark results and representative channel behavior. Section~\ref{sec:variants} summarizes exploratory variants, including longer training, manifold uncertainty, and the multi-scale TCN fork. Finally, Sections~\ref{sec:conclusion} and~\ref{sec:future} summarize the findings and outline future directions.

\section{Related Works}
\label{sec:related}

This section summarizes academic articles that discuss anomaly identification by evaluating space telemetry using the Soil Moisture Active Passive (SMAP) satellite and the Mars Science Laboratory (MSL) rover Dataset \cite{Hundman2018b}. 

The study by Hundman et al. demonstrates the application of Long Short-Term Memory (LSTM) networks in detecting spacecraft anomalies through using a combination of deep learning and nonparametric dynamic thresholding. Unlike conventional methods, this approach leverages the sequential nature of telemetry data to forecast upcoming values, identifying anomalies wherever the actual data significantly diverges from the forecast. The LSTM model captures long-range dependencies \cite{hua2019deep}, allowing it to learn intricate temporal relationships within the data. Dynamic thresholding adjusts to the normal behavior variations over time, enhancing the anomaly detection system's sensitivity and robustness. This dual-method strategy has shown improved precision in identifying both contextual and collective anomalies within multivariate telemetry data. Evaluations on the SMAP dataset validate the approach's efficacy, showcasing advances within the constraints of space missions and reinforcing its relevance to aerospace applications.

Recognizing the importance of long-range dependencies, recent research has also focused on improving model scalability to handle larger context windows. A notable example is the work by Hoh et al. \cite{hoh2025rethinking}, who propose a framework based on a Swin Transformer U-Net architecture (Swin1D-AD). By leveraging the hierarchical and shifted-window attention of the Swin Transformer \cite{he2022swin}, their model efficiently processes input windows of over 1000 timesteps, significantly larger than typical baselines. This enables the capture of distant temporal relationships within a reconstruction-based paradigm, demonstrating that larger contexts can improve detection performance. However, while this approach effectively expands the \textit{size}, it does not address the fundamental problem of the context's \textit{reliability}. As a reconstruction-based model, it remains vulnerable to context contamination; a sufficiently prolonged anomaly can still become the "new normal" that the model learns to reconstruct, leading to eventual score suppression.

Further advancing the state of the art, researchers have focused on integrating learned structural relationships with temporal dynamics to create richer representations of context. A prominent example is the Learning Graph Anomaly Transformer (LGAT) model by Wen et al. \cite{wen2025lgat}, which automatically learns a graph structure representing inter-variable dependencies and integrates it with an enhanced, U-Net-inspired Anomaly Transformer architecture. By simultaneously modeling these complex topological and temporal patterns, LGAT aims to build a more comprehensive representation of system state, thereby improving detection sensitivity based on reconstruction error and association discrepancies. Nevertheless, while this approach significantly enhances the modeling of intricate variable interactions, its performance is still predicated on the assumption that the immediate context window represents a normal operational state. This leaves it potentially vulnerable to the gradual score suppression characteristic of \textit{context contamination} during prolonged anomalies. Our work addresses a distinct but complementary challenge: how to model the context better and ensure the context's fundamental \textit{reliability}.

Neural Contextual Anomaly Detection (NCAD) is a cutting-edge framework for detecting anomalies in time series data by leveraging both supervised and unsupervised settings \cite{carmona2021neural}. NCAD distinguishes itself from traditional methods through its versatile training framework, which is underpinned by effective data augmentation heuristics. The architecture utilizes TCNs \cite{bai2018empirical} to derive latent representations from the temporal context and suspect windows. Anomalies are identified by comparing the embeddings of these windows and calculating the pseudo-Huber loss, which measures the distance-based anomaly score. This approach enables NCAD to leverage both labeled and unlabeled data, thereby enhancing its performance across multiple domains. Additionally, NCAD incorporates data augmentation techniques, such as Contextual Outlier Exposure (COE) and Point Outliers (PO), to inject synthetic anomalies, thereby enhancing the model's robustness. Empirical results display NCAD's superior accuracy and flexibility in detecting complex anomalies in real-world time series datasets, emphasizing its potential for broader applicability in anomaly detection tasks \cite{hendrycks2018deep}.

A significant research thrust addresses the growing computational complexity and parameter overhead of state-of-the-art models, which can be prohibitive in resource-constrained environments. A key contribution in this area is the Diminutive Memory-Enhanced Reconstruction (DiMER) model by Li et al. \cite{li2025contrast}, a lightweight framework designed specifically for efficiency. DiMER introduces a \textit{contrast memory} mechanism, which enables a compact encoder-decoder architecture to learn and retain distinct normal patterns effectively with minimal parameters. This memory is a core component of the reconstruction process, enhancing the model's ability to create a discriminative latent space by contrasting different normal sub-patterns. While this contrastive memory is highly effective for achieving state-of-the-art performance in a lightweight model, its primary function is to aid in efficient reconstruction; it does not explicitly address the failure mode where the model's input context window becomes progressively corrupted by a sustained anomaly.

To better capture dependencies across different temporal scales, hybrid architectures that combine attention and memory mechanisms have emerged as a powerful paradigm. The HYbrid Memory and Attention Network (HYMAN) framework, proposed by Li et al. \cite{li2025hyman}, exemplifies this strategy. It employs a local attention mechanism to model short-term patterns, while a parallelizable global memory module learns and stores key historical patterns to capture long-term dependencies. This design's simplified, auxiliary-free training process is a notable feature that improves engineering practicality. While this hybrid approach provides a powerful method for learning rich representations, the memory in HYMAN serves as an internal component to enrich the latent representation for reconstruction, rather than to validate the normalcy of the input context itself externally.

Finally, other hybrid models, such as AA-GAT, combine Graph Attention Networks (GAT) with an Autoformer to model spatial and temporal patterns, utilizing reconstruction error for scoring and contrastive learning as an auxiliary loss \cite{shang2024anomaly, ma2025time}. Its reliance on a predictive model that learns from a recent time window makes it similarly susceptible to context contamination during sustained faults.

\section{Dataset}
\label{sec:dataset}
\begin{figure}[t] % Try 't' (top) or 'b' (bottom) for better flow in two-column
    \centering
    \includegraphics[width=\linewidth]{telemetry.png}
    \caption{Visualization of SMAP Dataset features \cite{meng2019spacecraft}. Shows telemetry value (first feature) and one-hot encoded commands (remaining features).}
    \label{fig:smap_vis}
\end{figure}

The benchmark used in this study is composed of multiple telemetry channels from the public NASA SMAP/MSL collection \cite{Hundman2018b}, where each channel is treated as an independent dataset. Each channel file contains multivariate data, including the primary telemetry value and one-hot encoded commands (Figure~\ref{fig:smap_vis}). Our current approach focuses on extracting comprehensive temporal features from the primary univariate telemetry value. These channels exhibit two main anomaly types: contextual anomalies, where a point's deviation is only apparent given its surrounding context, and point anomalies, where a data point distinctly deviates from the bulk of the dataset regardless of context \cite{10400552}.

\section{Proposed Method}
\label{sec:proposed_method}

Our proposed methodology, termed Neural Contextual Anomaly Detection with Counterfactual Successor Memory (NCAD-CS), is a multi-stage methodology for anomaly detection in complex spacecraft telemetry. Building upon the contrastive-window formulation of the original NCAD framework \cite{carmona2021neural}, we retain a shared TCN encoder for full and context windows, but change the downstream question from context substitution to successor consistency. During training, normal contexts are stored together with the normal suspect segment that immediately followed them. During inference, a test context retrieves nearest normal contexts and their successors, producing a local counterfactual set of plausible normal futures. The observed suspect window is then scored by how far it deviates from this set, with complementary evidence from local raw-signal deviation and context distance. The overall workflow is depicted in Figure~\ref{fig:workflow}.

\begin{figure}[t]
    \centering
    \includegraphics[width=\linewidth]{mermaid-ai-1.png}
    \caption{Illustration of the algorithm workflow.}
    \label{fig:workflow}
\end{figure}

\subsection{Robust Multi-Domain Feature Engineering}
\label{subsec:feature_eng}

Recognizing that raw telemetry often lacks the expressive power to capture subtle deviations, we begin with a multi-domain feature engineering stage. This process transforms the raw univariate signal into a rich, high-dimensional representation, providing the model with a holistic view of the system's state across various temporal and spectral dimensions. Given a raw univariate telemetry channel $x_{\text{raw}} = \{x_1, x_2, \dots, x_T\}$, we extract a comprehensive set of features, including:
\begin{itemize}[noitemsep, topsep=0pt, leftmargin=*]
    \item \textit{Short-term rolling statistics} (e.g., mean, standard deviation, skew, kurtosis) to model immediate local signal characteristics.
    \item \textit{Long-term rolling statistics} over a significantly larger window to establish a stable baseline of the signal's intrinsic volatility \cite{zivot2003rolling}.
    \item \textit{Frequency-domain features} derived from both the Fast Fourier Transform (FFT) for global spectral properties and the Short-Time Fourier Transform (STFT) to identify localized changes in spectral energy and multimodal behavior \cite{chen2007feature}.
    \item \textit{Trend and complexity indicators}, such as rolling slope for trend analysis and wavelet energy ratios for capturing non-stationary, multi-scale patterns \cite{trokic2019wavelet}.
\end{itemize}

A critical aspect of this stage is ensuring numerical stability, particularly when dealing with real-world telemetry that may include periods of constant or near-constant values. Features relying on variance, such as the Z-score or higher-order moments like skewness, can be ill-defined for data with zero variance. Our implementation explicitly handles these edge cases by checking for near-zero variance before calculation and providing sensible defaults (e.g., zero skewness for a constant signal). This ensures the feature set remains robust and free of non-finite values that could disrupt model training.

This process yields a high-dimensional feature matrix, $F_{\text{all}} \in \mathbb{R}^{T \times D_{\text{full}}}$. Each feature column is then normalized via a Z-score transformation using statistics ($\bm{\mu}_{F}$, $\bm{\sigma}_{F}$) computed from the training data \cite{patro2015normalization}. The results are clipped to a stable range (e.g., [-10, 10]) to mitigate the dominance of outliers. Finally, for computational efficiency, variance-based feature selection reduces the dimensionality to the final input matrix $F_{\text{selected}} \in \mathbb{R}^{T \times k}$, where $k=64$ in our experiments.

\subsection{Windowing and Hybrid-Pooling TCN Encoder}
\label{subsec:enhanced_tcn_encoder}

The core of our anomaly detection system is a sophisticated encoder, $\phi(\cdot; \theta)$, designed to transform time series windows of fixed size into information-rich latent representations. Our encoder is built upon a Temporal Convolutional Network (TCN), an architecture well-suited for sequential data \cite{bai2018empirical}. Compared to recurrent architectures like LSTMs, TCNs offer several key advantages for this problem. Specifically, unlike the sequential, step-by-step processing inherent to LSTMs, TCNs apply 1D convolutions across the entire input sequence in parallel. This architectural difference significantly accelerates model training and inference, a critical requirement for handling the high-volume data streams common in space systems \cite{gopali2021comparison}. Furthermore, TCNs achieve a vast receptive field by stacking layers of dilated causal convolutions. This allows the model to capture extremely long-term dependencies—such as a slow system drift over thousands of timesteps—without the vanishing gradient problems that can still affect LSTMs over such long horizons \cite{chen2022vanishing}. The stable gradient flow and flexible receptive field make TCNs exceptionally well-suited for learning the complex, multi-scale temporal patterns characteristic of spacecraft telemetry. Our design, however, moves beyond a standard TCN to address the unique challenges of spacecraft telemetry, which exhibits a broad spectrum of signal dynamics, from quiescent periods to highly volatile events.

\subsection{Input Segmentation and Normalization}
The input to the encoder is a feature matrix, $F_{\text{selected}}$, which has been pre-processed and feature-engineered. This matrix is first segmented into overlapping windows of a fixed length, $L$. Each window, $w \in \mathbb{R}^{L \times k}$ (where $k$ is the number of features), is further subdivided into two distinct segments:
\begin{itemize}
    \item A context window, $w^{(c)} \in \mathbb{R}^{C \times k}$, representing the recent, presumably normal, history of the signal.
    \item A suspect window, $w^{(s)} \in \mathbb{R}^{S \times k}$, representing the immediate subsequent data points to be evaluated for anomalies.
\end{itemize}
Here, the total window length is $L = C + S$.

To ensure model stability and robustness, we integrate Layer Normalization \cite{ba2016layernormalization} directly within our TCN blocks. Unlike batch-dependent normalization schemes, Layer Normalization computes its statistics independently for each sample across its features. This makes the encoder's performance resilient to variations in batch composition and prevents issues with low-variance "flat-line" signals, common in telemetry data.

\begin{figure}[t] % Changed from h! to t for better placement
    \centering
    \includegraphics[width=0.9\linewidth]{TCN-Encoder.png}
    \caption{Conceptual Illustration of the Dual-Pathway Encoder Architecture.}
    \label{fig:encoder_architecture}
\end{figure}

\subsection{Pathway 1: The Hybrid Temporal Pathway}
A key innovation of our model lies in its dual-pathway architecture, which processes the intermediate feature maps from the TCN, denoted as $X_{\text{tcn}}$, through two concurrent and complementary pathways. This design is engineered to capture a multi-faceted understanding of the window's content, analyzing its temporal evolution and underlying statistical properties. The outputs of these pathways are then combined to form the final latent embedding.

This pathway extracts a rich, multi-resolution summary of the window's temporal dynamics. It recognizes that anomalies manifest in various ways over time—some are abrupt spikes, while others are gradual deviations. To capture this diversity, it employs a hybrid pooling strategy that combines three distinct views of the TCN output, $X_{\text{tcn}}$:

\begin{itemize}[noitemsep, topsep=0pt, leftmargin=*]
    \item \textit{Final State Representation:} The feature vector from the last time step of $X_{\text{tcn}}$ is extracted directly. Due to the TCN's causal structure and large receptive field, this vector is the model's most up-to-date representation after processing the entire sequence. It is therefore susceptible to the most recent changes and is critical for detecting sharp, transient point anomalies that global pooling operations might otherwise dilute.

    \item \textit{Dominant Feature Pooling (Max-Pooling):} Max-pooling is applied across the entire time axis of the feature maps. This operation identifies the most salient or dominant features present anywhere within the window, regardless of their exact temporal location. It effectively captures the peak intensity of events and provides a strong, time-invariant signal for persistent states \cite{gholamalinezhad2020pooling}.

    \item \textit{Global Summary Pooling (Average-Pooling):} Average-pooling is also applied across the time axis, providing a holistic summary of the window's overall pattern and central tendency. This helps identify long-term, slowly evolving states and provides a stable, baseline representation of the window's content.
\end{itemize}

The feature vectors resulting from these three pooling operations are immediately concatenated. This combined temporal vector is then projected through a dense layer with a non-linear activation function to produce a comprehensive temporal feature vector. This hybrid approach ensures the pathway is sensitive to transient and sustained phenomena.

\subsection{Pathway 2: The Statistical and Spectral Pathway}
This pathway operates parallel to the temporal path and captures the signal's underlying "texture" and phase-invariant characteristics. While the temporal pathway focuses on what happens over time, this pathway focuses on how the signal behaves statistically and in the frequency domain.

For each filter in the TCN output feature map, we compute a set of statistical moments and energy-like features over the time axis. This includes the mean, standard deviation, skewness, kurtosis, and spectral power. This process, inspired by the success of feature engineering in classical signal processing \cite{chaman2021truly}, provides a detailed signature of the signal's properties that is robust to minor temporal shifts or phase changes. For example, it can distinguish between a low-amplitude, high-frequency oscillation and a high-amplitude, low-frequency trend, even if their temporal averages are similar. The resulting statistical features are passed through their dedicated dense layer to form the spectral feature vector.

The feature vectors produced by the Hybrid Temporal Pathway and the Statistical \& Spectral Pathway are finally concatenated. This combined super-vector, which now contains a comprehensive description of the input window from multiple perspectives, is projected by a final bottleneck layer to yield the ultimate latent embedding, $z \in \mathbb{R}^{16}$.  

This dual-pathway, hybrid-pooling design is a cornerstone of our model's ability to detect a wide and diverse range of anomalies. By applying the same shared-weight encoder, $\phi$, to both the full window ($w$) and the context-only window ($w^{(c)}$), we obtain their respective latent embeddings, $z$ and $z^{(c)}$, enabling downstream comparison and anomaly scoring \cite{carmona2021neural}.

\subsection{Counterfactual Successor Memory and Scoring}
\label{subsec:successor_memory}
A critical failure mode in contrastive anomaly detection is \textit{context contamination}. As a sustained anomaly persists, it fills the context window, as illustrated in Figure~\ref{fig:context_suspect_contaminated}. A naive model comparing this contaminated context to the anomalous suspect window can perceive them as mutually consistent, causing the anomaly score to be suppressed and the fault to be missed.

NCAD-CS addresses this failure mode by replacing direct context substitution with Counterfactual Successor Memory (CSM). Instead of replacing the current context, CSM stores normal contexts together with the normal successor segments that followed them. At inference time, nearest normal contexts retrieve plausible normal futures, and the observed suspect segment is scored against those counterfactual successors. The process is realized through the following stages:

\begin{figure}[t]
    \centering
    \includegraphics[width=\linewidth]{figure_3.png}
    \caption{Illustration of a clean context window (yellow) preceding the suspect window (red) at the onset of a contextual anomaly.}
    \label{fig:context_suspect_clean}
\end{figure}

\begin{figure}[t]
    \centering
    \includegraphics[width=\linewidth]{figure_4.png}
    \caption{Illustration of a contaminated context window, which can mislead standard context-vs-suspect comparison methods.}
    \label{fig:context_suspect_contaminated}
\end{figure}

\textit{1) Successor Memory Construction:} Let each training window be decomposed into a context segment $w_i^{(c)} \in \mathbb{R}^{C \times k}$ and a successor segment $w_i^{(s)} \in \mathbb{R}^{S \times k}$. After contrastive encoder training, we encode each normal context as $z_i^{(c)}=\phi(w_i^{(c)};\theta)$ and store the pair $(z_i^{(c)}, w_i^{(s)})$ in memory. For computational efficiency, the implementation may subsample the memory to a maximum of 5000 windows per channel. A nearest-neighbor index is then fit over the context embeddings. The context reliability threshold $\tau_c$ is defined as the 99th percentile of leave-one-out nearest-neighbor distances among normal training contexts.

\textit{2) Counterfactual Successor Retrieval:} For a test window $w_t=(w_t^{(c)}, w_t^{(s)})$, the context embedding $z_t^{(c)}$ retrieves its $K$ nearest normal context embeddings. Let $\mathcal{N}_K(t)$ denote their indices and let \{w_j^{(s)}:j\in\mathcal{N}_K(t)\} be the retrieved normal successors. These successors form a local counterfactual set: plausible next segments under normal behavior for contexts similar to $w_t^{(c)}$. The implementation uses $K=8$ neighbors.

\textit{3) Successor Residual and Dispersion:} The primary successor residual is the minimum root-mean-square error between the observed successor and the retrieved normal successor set:
\begin{equation}
    r_{\text{succ}}(t) = \min_{j \\in \mathcal{N}_K(t)} \sqrt{\frac{1}{Sk}\left\|w_t^{(s)} - w_j^{(s)}\right\|_2^2}.
\end{equation}
We also compute the median retrieved successor, $\tilde{w}^{(s)}_t=\operatorname{median}_{j\in\mathcal{N}_K(t)} w_j^{(s)}$, and the dispersion of the retrieved successor set. The median successor residual is retained as a diagnostic, while the dispersion is used by the experimental manifold-uncertainty variant discussed in Section~\ref{sec:variants}.

\textit{4) Evidence Fusion:} Raw successor residuals are calibrated using robust statistics from the leave-one-out training memory distribution. Specifically, the successor evidence is a positive robust z-score, $z_{\text{succ}}$, computed from the median and interquartile range of normal calibration residuals. A complementary local evidence term, $z_{\text{local}}$, measures abrupt suspect deviations relative to the recent raw-signal context using robust local scale estimates. Finally, the context distance ratio
\begin{equation}
    \rho_c(t) = \frac{d_c(t)}{\tau_c}
\end{equation}
increases successor evidence when the current context is farther from the normal context manifold. The fused window score is
\begin{equation}
S_{\text{win}}(t)=\max\left[z_{\text{succ}}(t)g_c(t),\,0.8z_{\text{local}}(t)\right],
\end{equation}
where
\begin{equation}
g_c(t)=1+0.35\cdot\operatorname{clip}(\rho_c(t)-1,0,2).
\end{equation}
This fusion keeps the core successor-memory evidence while preserving sensitivity to abrupt point or mean-shift deviations that may be visible directly in the suspect segment.


\begin{figure}[t]
    \centering
    \includegraphics[width=\linewidth]{context_substitution_illustration.png}
    \caption{Conceptual view of counterfactual successor retrieval: similar normal contexts provide plausible normal successor segments against which the observed suspect segment is scored.}
    \label{fig:context_substitution}
\end{figure}

\subsection{Encoder Training}
\label{subsec:encoder_training}

The TCN encoder is trained using a contrastive learning objective \textit{before} Counterfactual Successor Memory is constructed. This training aims to sculpt a discriminative latent space where embeddings of normal operational sequences are clustered tightly, while those containing anomalies are pushed distinctly apart. 

To achieve this, we programmatically inject synthetic anomalies into the 16-sample suspect window of a random 70\% of the training windows in each batch. The injected anomalies are sampled from diverse types, including spikes, level shifts, variance changes, and stuck-values, with magnitudes scaled relative to the context window's standard deviation. This process creates a challenging set of paired examples:
\begin{itemize}[noitemsep, topsep=0pt, leftmargin=*]
    \item \textit{Positive Pairs ($y=0$):} A normal window $w_i$ and its corresponding context $w_i^{(c)}$. The model should learn to produce similar embeddings for these, i.e., $z_i \approx z_i^{(c)}$.
    \item \textit{Negative Pairs ($y=1$):} A synthetically modified window $w'_i$ (containing an anomaly) and its original context $w_i^{(c)}$. The model should learn to produce dissimilar embeddings, pushing $z'_i$ far from $z_i^{(c)}$.
\end{itemize}

We employ a contrastive loss function with a fixed margin, $m$, which formalizes this objective \cite{le2020contrastive}. For a given pair of embeddings, the loss $\mathcal{L}$ is defined as:
\begin{equation}
\mathcal{L}(z_1, z_2, y) = (1-y) \cdot \mathcal{D}(z_1, z_2)^2 + y \cdot \max(0, m - \mathcal{D}(z_1, z_2))^2
\label{eq:contrastive_loss}
\end{equation}
where $\mathcal{D}(\cdot, \cdot)$ is the distance metric (e.g., Euclidean distance) and $y$ is the binary label indicating if the pair is positive ($y=0$) or negative ($y=1$) \cite{kostadinov2015vector}.

This loss function intuitively penalizes two distinct behaviors. For positive pairs, the loss is simply the squared distance, driving the model to make the embeddings of normal windows and their contexts as close as possible. The loss is zero for negative pairs if the distance exceeds the margin $m$. If the distance is less than $m$, the loss term pushes them apart until they are at least a distance $m$ from each other. Training on diverse synthetic anomalies makes the encoder highly sensitive to deviations from established temporal patterns.

\subsection{Final Inference and Thresholding}
\label{subsec:inference_thresholding}
The final stage translates window-level successor-memory evidence into point-wise predictions through aggregation, smoothing, adaptive thresholding, and event filtering.

\textit{1) Aggregation and Smoothing:} Window scores $S_{\text{win}}$ are assigned to the suspect region of each window and averaged into a point-level score sequence $S_{\text{point}}$. The implementation then applies a short moving average with a 12-sample smoothing window. This smoothing is intentionally modest: it reduces isolated score jitter while preserving the onset of short telemetry events.

\textit{2) Training-Derived Event Threshold:} The memory construction stage produces leave-one-out calibration scores on normal training windows. These scores are passed through the same successor, local-deviation, and context-ratio fusion path used at inference. After aggregation and smoothing, the training event threshold $T_{\text{train}}$ is set to the 99th percentile of valid calibration scores. This gives a channel-specific normal-event threshold without using test labels.

\textit{3) Channel-Adaptive Elbow Score Floor:} A training-only threshold can be too low on channels whose normal successor residuals are nearly zero, causing broad false positives on flat or saturated channels. To prevent this, NCAD-CS computes an additional score floor from the geometry of the test-channel score distribution. The accepted implementation uses an adaptive elbow procedure over the sorted score curve and histogram separation. If the score distribution contains a clipped maximum-score plateau, the elbow is first estimated below the plateau so that saturated anomaly evidence is not hidden by an equal threshold. The final threshold is
\begin{equation}
    T_{\text{final}} = \max(T_{\text{train}}, T_{\text{elbow}}).
    \label{eq:t_final_revised}
\end{equation}
This design preserves the training-derived normality reference while avoiding a fixed global score-floor percentile.

\textit{4) Event-Level Filtering:} A time point is preliminarily flagged when $S_{\text{smooth},t} > T_{\text{final}}$. Contiguous flagged regions are retained if they satisfy a minimum run length of two points, exceed an extreme-score factor of 1.75, or contain sufficient score area above threshold. This event-level filter removes isolated noise while allowing high-magnitude point anomalies to pass.

This end-to-end process grounds final decisions in retrieved normal successors, local raw-signal evidence, and channel-adaptive threshold geometry.


\section{Experimentation}
\label{sec:exp}

This section outlines the experimental methodology and assesses the performance of the NCAD-CS approach in detecting various types of anomalies.

\subsection{Experimental Setup}
\label{subsec:exp_setup}
The experiments were conducted using telemetry channels from the SMAP and MSL dataset provided by NASA \cite{Hundman2018b}. Each channel was treated as an independent train-test problem. The accepted NCAD-CS configuration used the following settings:
\begin{itemize}[noitemsep, topsep=0pt, leftmargin=*]
    \item \textit{Windowing:} A full window size of 300 samples was used, consisting of a 284-sample context window and a 16-sample suspect/successor window. Windows were generated with a step size of 1.
    \item \textit{Feature Representation:} Each univariate telemetry channel was transformed into a 64-dimensional feature representation using the multi-domain feature extraction pipeline described in Section~\ref{subsec:feature_eng}.
    \item \textit{Encoder:} The reported default uses the hybrid-pooling TCN encoder with a 16-dimensional latent embedding, four TCN layers, 64 filters, kernel size 5, dropout 0.20, and Layer Normalization. A multi-scale TCN fork is evaluated separately in Section~\ref{sec:variants}.
    \item \textit{Training:} The encoder was trained with contrastive synthetic anomaly injection, batch size 32, AdamW optimization, and early stopping. The accepted full-channel result uses 15 training epochs; a 30-epoch run is reported as a stability check.
    \item \textit{Successor Memory:} Counterfactual Successor Memory used $K=8$ nearest normal contexts, a maximum of 5000 stored memory windows per channel, and a 99th-percentile context-distance threshold derived from leave-one-out normal-context distances.
    \item \textit{Scoring and Thresholding:} Successor residual evidence, local raw-signal deviation, and context distance were fused into window scores. Point scores were smoothed with a 12-sample moving average. The final event threshold used the 99th percentile of training calibration scores with the adaptive elbow score floor described in Section~\ref{subsec:inference_thresholding}. No channel labels were used for threshold selection.
\end{itemize}


\subsection{All-Channel Results and Representative Behavior}
\label{subsec:results_analysis}
We analyze the accepted 15-epoch Counterfactual Successor Memory configuration across the full SMAP/MSL channel set. Rather than selecting only success cases, we use representative channels to show both the strengths and the remaining failure modes of the pipeline. Excluding one channel without valid labels, the run evaluated 81 telemetry channels and produced the aggregate status shown in Table~\ref{tab:current_aggregate_status}.

\begin{table}[h!]
\centering
\caption{Full-channel evaluation status for the accepted 15-epoch NCAD-CS Counterfactual Successor Memory configuration.}
\label{tab:current_aggregate_status}
\begin{tabular}{@{}lr@{}}
\toprule
Metric & Value \\
\midrule
Evaluated channels & 81 \\
Macro F1 & 0.3883 \\
Median F1 & 0.2875 \\
Micro precision & 0.3930 \\
Micro recall & 0.4587 \\
Micro F1 & 0.4233 \\
Channels with F1 $\\geq$ 0.5 & 34 \\
Channels with F1 $\\geq$ 0.8 & 16 \\
Zero-F1 channels & 12 \\
\bottomrule
\end{tabular}
\end{table}

The strongest channels in this configuration include \textit{B-1}, \textit{A-7}, \textit{D-4}, \textit{D-1}, \textit{D-3}, \textit{P-3}, \textit{F-5}, \textit{P-7}, \textit{M-7}, \textit{E-11}, \textit{E-7}, and \textit{D-8}. These results show that successor retrieval can be highly effective when similar normal contexts imply a coherent set of normal successors. However, the aggregate metrics also show that the current method is not yet robust across the full benchmark, with 12 zero-F1 channels and several low-precision over-firing cases.

... (file continues; backup contains full original)
